\section{Introduction}\label{introduction}

Process owners rarely get to test changes risk-free. Before changing
staffing rules, routing policies, automation levels, or escalation
procedures, they need a way to estimate what the intervention might do.
Business process simulation supports this by generating hypothetical
executions of a process under explicit assumptions. In process mining,
simulation is attractive because historical event logs already contain
activities, resources, timestamps, case arrivals, and observed variants.
A data-driven simulator can be calibrated from the same logs used for
process discovery and conformance checking.

Despite this empirical grounding, many BPS approaches remain more
process-centric than resource-centric. They model a process as a
control-flow structure with timing and resource parameters attached to
it. This is useful for many what-if analyses, but it can hide the
decentralized character of organizational work. Individual resources
differ in capabilities, workload, availability, handover behavior, and
local routines. A case does not move only because a global model marks
the next activity as enabled. It moves because a person, team, or
software resource accepts work, prioritizes it, performs it, and passes
it to another actor.

Recent agent-based approaches respond to this limitation by representing
resources as agents. AgentSimulator, for example, moves data-driven BPS
toward an agent-based view by learning agent behavior from event logs
\cite{kirchdorfer2024agentsimulator}. In parallel, work on
LLM-empowered agent-based modeling argues that LLMs can provide richer
agent cognition, communication, and decision-making in simulations
\cite{gao2024llmabm}. Together, these lines of work suggest a concrete
research opportunity: a business process can be simulated as a set of
log-grounded resource agents, with LLM-style modules used for local
decisions and handovers rather than for free-form process generation.

The opportunity also creates a methodological risk. If LLM agents freely
invent process behavior, the resulting simulation may be fluent but not
reproducible, valid, or faithful to the observed process. For
process-mining applications, every simulated action must remain
compatible with an event-log schema, a feasible activity-resource
assignment, and a measurable quality framework. I therefore treat LLM
agents as constrained decision modules. They operate inside action
spaces derived from event logs, and their generated event logs remain
subject to the same log-distance measures used for conventional BPS.

The research question is:

\textbf{How can log-derived resource-agent profiles and local behavior
priors be used to constrain LLM-augmented agent decisions so that a
resource-centric multi-agent simulator generates
process-mining-compatible event logs and interpretable reasoning and
handover traces?}

The contribution is a design-science artifact and an initial controlled
evaluation. The artifact, denoted as \(H_{R-LLM}\),
represents resources as agents with learned capabilities, service-time
distributions, local resource priors, handover priors, and memory. It
compares three executable policies: a centralized baseline, a
log-derived agent-profile policy, and a constrained LLM-agent proxy. The
LLM-agent proxy is a deterministic, inspectable stand-in for the
decision interface that a real LLM would occupy. That choice keeps the
first evaluation reproducible while testing whether an LLM-style local
decision module can enter a BPS pipeline without breaking event-log
compatibility.

The evaluation uses two public event logs: AcademicCredentials from the
artifact accompanying Chapela-Campa et al.'s BPS quality framework
\cite{chapela2025quality}, and LoanApp from AgentSimulator
\cite{kirchdorfer2024agentsimulator}. The experiments compare simulated
logs against held-out logs with lightweight prototype metrics and formal
BPS log-distance measures, then add what-if stress tests to examine
scenario response. The LLM-agent proxy does not dominate the baselines.
It is competitive on control-flow reproduction and obtains the best
workforce distance on both datasets, while the agent-profile and central
baselines remain stronger on several timing and control-flow measures.
Under some LoanApp high-load conditions the proxy responds well, but
that advantage does not survive every stress scenario. The mixed result
is the point. LLM-agent process simulation needs both BPS log-quality
metrics and LLM-agent behavioral evaluation, because its value lies in
log reproduction, scenario response, reasoning traces, and handover
observability.
